\begin{center}
  {\Large\itshape Alban « la féodale » -- XIII\textsuperscript{ème} siècle}
\end{center}

\vspace{1.5em}

\noindent
\lettrine{A}{u} civil -- comme Nougaret -- Alban est englobé dans le Comté
de Castres. Il se pourrait qu'à la suite de la croisade des Albigeois,
Philippe de Montfort, neveu du trop fameux Simon, ait donné en 1229,
à la mort de son oncle, la Seigneurie d'Alban à l'Abbaye Saint Benoît
de Castres. C'est ce début du XIII\textsuperscript{ème} siècle que les
vicomtes d'Albi et d'Ambialet accordent des franchises aux villes de
leur ressort. Alban a dû posséder, dès cette époque, une ébauche
d'organisation municipale en la personne de deux consuls. Au
XVII\textsuperscript{ème} siècle, ces magistrats sont choisis par le juge
mage de la sénéchaussée de Castres sur une liste de quatre candidats
présentés par les électeurs d'Alban.

\vspace{0.8em}

\noindent
\lettrine{L}{a} justice est rendue par l'Abbé de Castres. Mais, au
XIV\textsuperscript{ème} siècle, il y a concurremment le sénéchal qui
l'exerce. Assez souvent, les ressortissants préféreront la justice
d'Eglise.

\vspace{0.8em}

\noindent
Deux bulles de Clément IV en 1266 et 1267 confirmant l'accord de 1231
entre l'évêque d'Albi et l'Abbé de Castres stipulent que St Michel de
Nougaret comme Ste Marie d'Alban demeurent sous juridiction
bénédictine. En 1274, Pierre Raymond de Janes y exerce la fonction
de notaire. En juin 1286, un grand seigneur ecclésiastique,
l'archevêque de Bourges métropolitain d'Albi passe par Alban, venant
de Montfranc où il a dit la messe chez les Hospitaliers à St Léonard
du Carmé. C'est avec une forte escorte qu'il traverse Alban pour se
rendre à Ambialet.
